\section{Reduction to the reciprocal family}
\label{sec:reduction}

Here I show that a globally periodic equation \eqref{eq:family} is the shift
or of the reciprocal form \eqref{eq:familya} with palindromic coefficients,
and that it is enough to treat supports with $\gcd(S\cup\{k\})=1$
(\autoref{cor:red-primitive}).

Most of this is known. The reduction to the reciprocal form and palindromy
are the Reduction Theorem (Theorem~2.4) of \citet{CimaGasullManosas2004}, and
\citet[Theorem~4]{CsornyeiLaczkovich2001} proved $|c_j|=|c_{k-j}|$ for
complex coefficients. I follow the route of Cima, Gasull and Ma\~nosas
(preimages of a one-parameter family of states, a Lyapunov function, and the
derivative at a fixed point), and give self-contained proofs because the
linearisation is needed again in \autoref{sec:small} and palindromy follows
from a short period-sum argument. That the linearisation of a periodic
recurrence at a fixed point has finite order goes back to
\citet[Theorem~3.2]{KurshanGopinath1974}; it was used by
\citet{KocicLadasRodrigues1993}, \citet[Lemma~2]{CsornyeiLaczkovich2001} and,
with affine normal forms at order two, \citet{RubioMasseguManosa2007}.

\subsection{The state map}

Throughout, the coefficients of \eqref{eq:family} are nonnegative, with
$A_0+\sum_{i\ge1}A_i>0$ and $B_0+\sum_{i\ge1}B_i>0$. Note that $A_{i+1}$ and
$B_{i+1}$ are the coefficients of $x_{n+i}$, for $0\le i\le k-1$; in particular
$A_1$ and $B_1$ multiply the oldest term $x_n$. Write
$\Omega=(0,\infty)^k$, and for $u=(u_1,\dots,u_k)\in\Omega$ put
\begin{equation}
g(u)=\frac{A_0+\sum_{i=1}^kA_iu_i}{B_0+\sum_{i=1}^kB_iu_i},
\qquad
F(u)=\bigl(u_2,\dots,u_k,\,g(u)\bigr).
\end{equation}
For $u\in\Omega$ the numerator and denominator of $g(u)$ are both positive,
since each contains a positive coefficient multiplied by a positive number. So
$F$ maps $\Omega$ to itself, every positive initial window
$(x_0,\dots,x_{k-1})$ extends to a unique positive solution of
\eqref{eq:family}, and the window $(x_n,\dots,x_{n+k-1})$ of that solution is
$F^{\,n}(x_0,\dots,x_{k-1})$. The equation is \textit{globally $p$-periodic}
($p\ge1$) when every positive solution satisfies $x_{n+p}=x_n$ for all $n$, and
this holds precisely when $F^{\,p}=\id$ on $\Omega$. In particular a globally
periodic $F$ maps $\Omega$ onto itself. The first step of the reduction,
\autoref{thm:red-dichotomy}, uses only this surjectivity; the second,
\autoref{thm:red-shift}, uses periodicity itself.

\subsection{The dichotomy}

The idea is to ask for a preimage of a single one-parameter family of states,
$(1,\dots,1,t)$ with $t>0$. That alone forces most of the coefficients to
vanish.

\begin{theorem}[dichotomy]
\label{thm:red-dichotomy}
Suppose $F(\Omega)=\Omega$. Then one of the following holds:
\begin{enumerate}
\item[(a)] $A_1=0$, $B_0=0$ and $B_i=0$ for all $i\ge2$, so that $B_1>0$ and
\eqref{eq:family} reads
$B_1x_{n+k}x_n=A_0+A_2x_{n+1}+\cdots+A_kx_{n+k-1}$;
\item[(b)] $B_1=0$, $A_0=0$ and $A_i=0$ for all $i\ge2$, so that $A_1>0$ and
\eqref{eq:family} reads
$x_{n+k}=A_1x_n/\bigl(B_0+B_2x_{n+1}+\cdots+B_kx_{n+k-1}\bigr)$.
\end{enumerate}
\end{theorem}

\begin{proof}
Fix $t>0$. A preimage of $(1,\dots,1,t)$ under $F$ must have the form
$(u,1,\dots,1)$ with $u>0$ and $g(u,1,\dots,1)=t$. Writing
$A'=\sum_{i\ge2}A_i$ and $B'=\sum_{i\ge2}B_i$ (both zero when $k=1$), the
hypothesis therefore says that for every $t>0$ there is a $u>0$ with
\begin{equation}
t\,(B_0+B_1u+B')=A_0+A_1u+A'.
\label{eq:red-dagger}
\end{equation}
There are four cases, according to the signs of $A_1$ and $B_1$.

If $A_1=B_1=0$, then $u$ drops out of \eqref{eq:red-dagger}, and taking $t=1$
and $t=2$ gives $B_0+B'=A_0+A'=2(B_0+B')$. Hence $B_0+B'=0$, so all
the $B_i$ vanish, which is excluded.

If $A_1=0<B_1$, then $tB_1u=A_0+A'-t(B_0+B')$ with $u>0$, so
$t(B_0+B')<A_0+A'$ for every $t>0$. Letting $t\to\infty$ gives
$B_0+B'\le0$, that is $B_0=0$ and $B_i=0$ for $i\ge2$: this is case (a).

If $B_1=0<A_1$, then symmetrically $A_1u=t(B_0+B')-A_0-A'$ with $u>0$, so
$A_0+A'<t(B_0+B')$ for every $t>0$. Letting $t\to0^+$ gives $A_0+A'\le0$,
that is $A_0=0$ and $A_i=0$ for $i\ge2$: this is case (b).

Finally, suppose $A_1>0$ and $B_1>0$, and put $\lambda=A_1/B_1$. At $t=\lambda$
the terms in $u$ cancel from \eqref{eq:red-dagger}, leaving
$\lambda(B_0+B')=A_0+A'$. Subtracting this from \eqref{eq:red-dagger} at
$t=\lambda+1$ leaves $(B_0+B')+B_1u=0$ for some $u>0$, which is impossible.

In case (a) the denominator is $B_1x_n$, so $B_1>0$; in case (b) the numerator
is $A_1x_n$, so $A_1>0$.
\end{proof}

\subsection{Branch (b) is the shift}

Genuine periodicity, rather than surjectivity, now removes the rest of branch
(b). Two simple tools are used. The first is a Lyapunov function: the
coordinate sum $\Sigma(u)=\sum_iu_i$ changes under one step by
\begin{equation}
\Sigma(F(u))-\Sigma(u)=g(u)-u_1 ,
\label{eq:red-lyap}
\end{equation}
so if $g(u)<u_1$ for every $u\in\Omega$ then $\Sigma(F^{\,p}(u))<\Sigma(u)$ for
every $p\ge1$, and $F^{\,p}=\id$ is impossible; the same holds if $g(u)>u_1$
throughout. The second is the derivative at an interior fixed point: if
$F^{\,p}=\id$ and $F(u^*)=u^*$, the chain rule gives $DF(u^*)^p=I$.

\begin{theorem}[branch (b) is the shift]
\label{thm:red-shift}
A globally periodic equation of type (b) has $A_1=B_0$ and $B_i=0$ for all
$i\ge2$; that is, it is the shift $x_{n+k}=x_n$.
\end{theorem}

\begin{proof}
Let the equation be globally $p$-periodic, so that $F^{\,p}=\id$ on $\Omega$.
In type (b), $g(u)=A_1u_1/\bigl(B_0+\sum_{i\ge2}B_iu_i\bigr)$. Write
$\beta=\sum_{i\ge2}B_i$.

Suppose first that $\beta=0$. Then $B_0>0$ (the denominator is not zero) and
$g(u)=(A_1/B_0)u_1$. If $A_1<B_0$ then $g(u)<u_1$ throughout $\Omega$, and if
$A_1>B_0$ then $g(u)>u_1$ throughout; both are excluded by
\eqref{eq:red-lyap}. So $A_1=B_0$, which is the conclusion.

Suppose next that $\beta>0$ and $A_1\le B_0$. Then for $u\in\Omega$ the
denominator exceeds $B_0\ge A_1$, so $g(u)<u_1$ throughout $\Omega$, which is
again excluded.

There remains the case $\beta>0$, $A_1>B_0$, which must be shown impossible.
Here $x^*=(A_1-B_0)/\beta>0$ satisfies $B_0+\beta x^*=A_1$, so the point
$u^*=(x^*,\dots,x^*)$ is fixed by $F$. The map $F$ is rational, with a
denominator that is positive on the open set $\Omega$, so it is smooth there,
and $F^{\,p}=\id$ gives $DF(u^*)^p=D(F^{\,p})(u^*)=I$ by the chain
rule. At $u^*$ the denominator of $g$ equals $A_1$, so
$\partial g/\partial u_1=1$ and $\partial g/\partial u_i=-x^*B_i/A_1$ for
$i\ge2$; thus, with $d_i=B_ix^*/A_1\ge0$,
\[
DF(u^*)\,w=\Bigl(w_2,\dots,w_k,\;w_1-\sum_{i\ge2}d_iw_i\Bigr).
\]
Put $\phi(\lambda)=\lambda^k+\sum_{i\ge2}d_i\lambda^{i-1}$. Then $\phi(0)=0$
and $\phi(1)=1+\beta x^*/A_1>1$, so by the intermediate value theorem there is
a $\lambda\in(0,1)$ with $\phi(\lambda)=1$. The vector
$v=(1,\lambda,\dots,\lambda^{k-1})$ then satisfies $DF(u^*)\,v=\lambda v$,
since the last coordinate of $DF(u^*)\,v$ is
$1-\sum_{i\ge2}d_i\lambda^{i-1}=\lambda^k$. Hence $v=DF(u^*)^pv=\lambda^pv$,
so $\lambda^p=1$, which is impossible for $0<\lambda<1$.
\end{proof}

\subsection{The reciprocal form and its normalisation}

\begin{theorem}[reduction]
\label{thm:reduction}
Every globally periodic equation of the form \eqref{eq:family} is either the
shift $x_{n+k}=x_n$ (that is, $A_1=B_0>0$ and every other coefficient is zero),
or satisfies $A_1=B_0=0$ and $B_i=0$ for $i\ge2$, with $B_1>0$. In the second
case, dividing through by $B_1$, its positive solutions are exactly the positive
solutions of
\begin{equation}
x_{n+k}\,x_n=c_0+\sum_{j\in S}c_j\,x_{n+j},
\qquad c_j>0 \text{ for } j\in S\subseteq\{1,\dots,k-1\},\quad c_0\ge0,
\label{eq:familya}
\end{equation}
where
\begin{equation}
c_0=\frac{A_0}{B_1},\qquad
c_j=\frac{A_{j+1}}{B_1}\quad(1\le j\le k-1),\qquad
S=\{\,j\in\{1,\dots,k-1\} : A_{j+1}>0\,\},
\label{eq:red-normalise}
\end{equation}
and $S\ne\emptyset$ or $c_0>0$.
\end{theorem}

\begin{proof}
A globally periodic $F$ maps $\Omega$ onto itself, so \autoref{thm:red-dichotomy}
applies. In case (b), \autoref{thm:red-shift} gives the shift. In case (a) the
equation is $B_1x_{n+k}x_n=A_0+\sum_{j=1}^{k-1}A_{j+1}x_{n+j}$, and dividing by
$B_1>0$ gives \eqref{eq:familya} with the data \eqref{eq:red-normalise}, after
dropping the terms with $A_{j+1}=0$. The last condition is the hypothesis that
the numerator of \eqref{eq:family} is not identically zero (recall $A_1=0$).
\end{proof}

From now on an equation of the form \eqref{eq:familya} is described by its
\textit{normalised data}: the order $k$, the support $S$ and the coefficients
$c_0$ and $c_j$ ($j\in S$); it is convenient to put $c_j=0$ for
$j\in\{1,\dots,k-1\}\setminus S$, and $\theta\in\{0,1\}$ records whether
$c_0>0$, as in the introduction. The
normalisation \eqref{eq:red-normalise} is a division of numerator and
denominator by $B_1$, not a change of variable, so the normalised equation has
the same solutions as the original one, and global periodicity of the one is
global periodicity of the other. When $k=1$ the support is empty and
\eqref{eq:familya} is $x_{n+1}x_n=c_0$.

\subsection{The linearisation}

The reciprocal family always has an interior equilibrium. Write
$s=\sum_{j\in S}c_j$. Since $s+c_0>0$, the equation $x^2=c_0+s x$
has exactly one positive root,
\begin{equation}
x^*=\tfrac12\Bigl(s+\sqrt{s^2+4c_0}\Bigr),
\end{equation}
and the constant sequence $x_n=x^*$ is a solution of \eqref{eq:familya}. Put
\begin{equation}
\eps_j=\frac{c_j}{x^*}\quad(1\le j\le k-1),
\qquad\text{so that}\qquad
\sum_{j=1}^{k-1}\eps_j=\frac{s}{x^*}=1-\frac{c_0}{x^{*2}}\le1,
\label{eq:red-ej}
\end{equation}
with equality exactly when $c_0=0$. The \textit{linearisation} of
\eqref{eq:familya} at $x^*$ is the linear recurrence
\begin{equation}
y_{n+k}+y_n=\sum_{j=1}^{k-1}\eps_j\,y_{n+j},
\label{eq:red-linrec}
\end{equation}
whose characteristic polynomial is
\begin{equation}
\chi(\lambda)=\lambda^k+1-\sum_{j=1}^{k-1}\eps_j\,\lambda^j .
\label{eq:red-chi}
\end{equation}
Every initial window $(y_0,\dots,y_{k-1})$ extends to a unique solution of
\eqref{eq:red-linrec}, and the window advances by the companion matrix $\Gamma$,
\[
\Gamma(w_1,\dots,w_k)=\Bigl(w_2,\dots,w_k,\;-w_1+\sum_{j=1}^{k-1}\eps_jw_{j+1}\Bigr).
\]

\begin{theorem}[linearisation]
\label{thm:spectral}
Suppose \eqref{eq:familya} is globally $p$-periodic. Then:
\begin{enumerate}
\item[(i)] $\Gamma^{\,p}=I$, so every solution of \eqref{eq:red-linrec} is
$p$-periodic;
\item[(ii)] every complex root $\lambda$ of $\chi$ satisfies $\lambda^p=1$.
\end{enumerate}
\end{theorem}

\begin{proof}
(i) By \autoref{thm:reduction} the state map of \eqref{eq:familya} is
$F(u)=\bigl(u_2,\dots,u_k,\,g(u)\bigr)$ with
$g(u)=\bigl(c_0+\sum_jc_ju_{j+1}\bigr)/u_1$, a rational map whose denominator
is positive on $\Omega$, and it fixes $u^*=(x^*,\dots,x^*)$. Using
$x^{*2}=c_0+\sum_jc_jx^*$, the partial derivatives of $g$ at $u^*$ are
$\partial g/\partial u_1=-x^{*2}/x^{*2}=-1$ and
$\partial g/\partial u_{j+1}=c_j/x^*=\eps_j$, so $DF(u^*)=\Gamma$. Since
$F^{\,p}=\id$ on the open set $\Omega$, the chain rule gives
$\Gamma^{\,p}=D(F^{\,p})(u^*)=I$.

(ii) If $\chi(\lambda)=0$, then $v=(1,\lambda,\dots,\lambda^{k-1})$ satisfies
$\Gamma v=\lambda v$, since the last coordinate of $\Gamma v$ is
$-1+\sum_j\eps_j\lambda^j=\lambda^k$. So $v=\Gamma^{\,p}v=\lambda^pv$, and
$\lambda^p=1$.
\end{proof}

\subsection{Palindromy}

The first structural consequence of \autoref{thm:spectral} is that the
coefficients are palindromic. The reason is that a linear map of finite order
preserves a positive definite quadratic form, so its inverse is its adjoint.
Here the form can be written down explicitly: on $p$-periodic sequences the
\textit{period sum}
\[
\langle u,v\rangle=\sum_{n=0}^{p-1}u_nv_n
\]
is positive definite, and it is unchanged by shifting both sequences,
$\sum_{n<p}u_{n+a}v_{n+a}=\langle u,v\rangle$, since $n\mapsto u_nv_n$ is
$p$-periodic. The recurrence \eqref{eq:red-linrec} must therefore be the same when
it is run backwards, and that is palindromy.

\begin{theorem}[palindromy]
\label{thm:palindromy}
If \eqref{eq:familya} is globally periodic, then $c_j=c_{k-j}$ for
$1\le j\le k-1$. In particular the support is palindromic, $S=k-S$; in the
notation of \eqref{eq:family}, $A_i=A_{k+2-i}$ for $2\le i\le k$.
\end{theorem}

\begin{proof}
Let the equation be globally $p$-periodic. By \autoref{thm:spectral}(i), every
real solution of \eqref{eq:red-linrec} is $p$-periodic. For a solution $v$, put
\[
\rho_n=v_n+v_{n+k}-\sum_{j=1}^{k-1}\eps_j\,v_{n+k-j}.
\]
Let $u$ be another solution, and expand $\sum_{n<p}u_n\rho_n$ into three period
sums. Shifting the first by $k$ and the terms of the third by $j$,
\[
\sum_{n<p}u_nv_n=\sum_{n<p}u_{n+k}v_{n+k},\qquad
\sum_{n<p}u_nv_{n+k-j}=\sum_{n<p}u_{n+j}v_{n+k},
\]
so that
\[
\sum_{n<p}u_n\rho_n
=\sum_{n<p}v_{n+k}\Bigl(u_{n+k}+u_n-\sum_{j=1}^{k-1}\eps_ju_{n+j}\Bigr)=0,
\]
because $u$ solves \eqref{eq:red-linrec}. On the other hand, using
\eqref{eq:red-linrec} for $v$ to replace $v_n+v_{n+k}$, and reindexing the sum
$j\mapsto k-j$,
\[
\rho_n=\sum_{j=1}^{k-1}\bigl(\eps_j-\eps_{k-j}\bigr)\,v_{n+j}.
\]
This is a linear combination of shifts of $v$, so $\rho$ is itself a solution of
\eqref{eq:red-linrec}. Taking $u=\rho$ gives $\sum_{n<p}\rho_n^2=0$, so $\rho_n=0$ for
$n<p$, and hence for all $n$ by periodicity.

Now fix $1\le j\le k-1$ and let $v$ be the solution with initial window
$v_j=1$ and $v_i=0$ for the other $i<k$. Then
$0=\rho_0=\eps_j-\eps_{k-j}$, and multiplying by $x^*$ gives $c_j=c_{k-j}$. The
statement about $S$ follows because $S=\{j:c_j>0\}$, and the statement about the
$A_i$ from \eqref{eq:red-normalise}.
\end{proof}

\subsection{Known types and primitive supports}

The five equations \eqref{eq:fivefamilies} have a simple description in terms
of normalised data. Recall that the \textit{$\ell$-fold unfolding} of an equation
$x_{n+k'}=f(x_n,x_{n+1},\dots,x_{n+k'-1})$ of order $k'$ is the equation
$x_{n+\ell k'}=f(x_n,x_{n+\ell},\dots,x_{n+(k'-1)\ell})$ of order $\ell k'$; its solutions
are exactly the sequences whose $\ell$ interleaved subsequences
$m\mapsto x_{m\ell+r}$ ($0\le r<\ell$) each solve the original equation.
(\citet{Gasull2021} describes this operation without a name, and uses the
word \emph{unfold} for the rescaled family $y_n=\alpha x_n$ instead.)

\begin{definition}
\label{def:red-known}
The normalised data $(k,S,c)$ of an equation \eqref{eq:familya} is of
\textit{known type} if one of the following holds:
\begin{enumerate}
\item[(K1)] $S=\emptyset$;
\item[(K2)] $k=2m$, $S=\{m\}$, and $c_0=0$ or $c_0=c_m^2$;
\item[(K3)] $k=3m$, $S=\{m,2m\}$, $c_m=c_{2m}$ and $c_0=c_m^2$.
\end{enumerate}
\end{definition}

\begin{lemma}
\label{lem:red-members}
If the normalised data of an equation of type (a) is of known type, then the
equation is equivalent, under a rescaling $y_n=x_n/\lambda$ ($\lambda>0$) and
an unfolding, to one of the equations \eqref{eq:fivefamilies}. In terms of the
coefficients of \eqref{eq:family}, with $m$ as in \autoref{def:red-known}:
in case (K1), $A_i=0$ for all $i\ge1$; in case (K2), $A_i=0$ for
$i\notin\{0,m+1\}$, and $A_0=0$ or $A_0B_1=A_{m+1}^2$; in case (K3), $A_i=0$
for $i\notin\{0,m+1,2m+1\}$, $A_{m+1}=A_{2m+1}$ and $A_0B_1=A_{m+1}^2$.
\end{lemma}

\begin{proof}
If $x$ solves \eqref{eq:familya}, then $y_n=x_n/\lambda$ solves
$y_{n+k}y_n=c_0/\lambda^2+\sum_{j\in S}(c_j/\lambda)\,y_{n+j}$. In case (K1),
$c_0>0$ by \autoref{thm:reduction}; with $\lambda=\sqrt{c_0}$ the equation
becomes $y_{n+k}=1/y_n$, the $k$-fold unfolding of $y_{n+1}=1/y_n$. In case
(K2), $\lambda=c_m$ gives $y_{n+2m}y_n=y_{n+m}$ or $y_{n+2m}y_n=1+y_{n+m}$,
the $m$-fold unfoldings of $y_{n+2}=y_{n+1}/y_n$ and of the Lyness equation
$y_{n+2}=(1+y_{n+1})/y_n$. In case (K3), $\lambda=c_m$ gives
$y_{n+3m}y_n=1+y_{n+m}+y_{n+2m}$, the $m$-fold unfolding of
$y_{n+3}=(1+y_{n+1}+y_{n+2})/y_n$. The coefficient conditions follow from
\eqref{eq:red-normalise}: $c_0=c_m^2$ reads $A_0/B_1=A_{m+1}^2/B_1^2$.
\end{proof}

The last step of the reduction removes common factors from the support. If
every element of $S$ and $k$ itself are multiples of some $\ell>1$, then the
equation only ever couples terms whose indices agree modulo $\ell$, and it splits
into $\ell$ independent copies of an equation of order $k/\ell$, as
\citet{CimaGasullManosas2004} observe at order four, where they split a
solution into its odd and even subsequences. Call $S$
\textit{primitive} (at order $k$) when $\gcd(S\cup\{k\})=1$.

\begin{lemma}[folding]
\label{lem:unfold}
Let \eqref{eq:familya} have normalised data $(k,S,c)$, and let
$\ell=\gcd(S\cup\{k\})$ (so $\ell=k$ when $S=\emptyset$). Put $k'=k/\ell$,
$S'=\{j/\ell:j\in S\}$, $c'_0=c_0$ and $c'_{j'}=c_{j'\ell}$ for $j'\in S'$. Then:
\begin{enumerate}
\item[(i)] equation \eqref{eq:familya} is the $\ell$-fold unfolding of the
equation
\begin{equation}
x_{n+k'}\,x_n=c'_0+\sum_{j'\in S'}c'_{j'}\,x_{n+j'}
\label{eq:red-folded}
\end{equation}
of order $k'$, which has the same $\theta$;
\item[(ii)] if \eqref{eq:familya} is globally $p$-periodic, then so is
\eqref{eq:red-folded};
\item[(iii)] $S'$ is primitive at order $k'$, and if $c_j=c_{k-j}$ for all $j$,
then $c'_{j'}=c'_{k'-j'}$ for all $j'$ (in particular $S'$ is palindromic);
\item[(iv)] if the data $(k',S',c')$ is of known type, then so is $(k,S,c)$.
\end{enumerate}
\end{lemma}

\begin{proof}
(i) Every element of $S$ is a multiple of $\ell$, and so is $k$. Writing
$n=m\ell+r$ with $0\le r<\ell$ and $j=j'\ell$, equation \eqref{eq:familya} at $n$ reads
\begin{equation*}
x_{(m+k')\ell+r}\,x_{m\ell+r}=c'_0+\sum_{j'\in S'}c'_{j'}\,x_{(m+j')\ell+r},
\end{equation*}
which is \eqref{eq:red-folded} at $m$ for the subsequence $m\mapsto x_{m\ell+r}$.

(ii) Let $z$ be a positive solution of \eqref{eq:red-folded}, and put
$x_n=z_{\lfloor n/\ell\rfloor}$. Each subsequence $m\mapsto x_{m\ell+r}$ equals $z$,
so by (i) $x$ is a positive solution of \eqref{eq:familya}, and it is therefore
$p$-periodic, hence also $p\ell$-periodic. So
$z_{m+p}=x_{m\ell+p\ell}=x_{m\ell}=z_m$ for every $m$.

(iii) and (iv) are immediate from the definitions:
$\gcd(S'\cup\{k'\})=\gcd(S\cup\{k\})/\ell=1$, and for (iv) one replaces $m$
by $m\ell$ in \autoref{def:red-known}.
\end{proof}

\begin{corollary}
\label{cor:red-primitive}
Suppose that every globally periodic equation \eqref{eq:familya} whose support
$S$ is primitive and whose coefficients satisfy $c_j=c_{k-j}$ for all $j$ has
normalised data of known type. Then every globally periodic equation of the
form \eqref{eq:family} is either the shift $x_{n+k}=x_n$, or is equivalent,
under a rescaling and an unfolding, to one of the equations
\eqref{eq:fivefamilies}, with coefficients as in \autoref{lem:red-members}.
\end{corollary}

\begin{proof}
Let \eqref{eq:family} be globally $p$-periodic. By \autoref{thm:reduction} it
is either the shift or has the normalised form \eqref{eq:familya}, which is
globally $p$-periodic. In the second case, \autoref{thm:palindromy} gives
$c_j=c_{k-j}$, and \autoref{lem:unfold} gives a globally $p$-periodic folded
equation \eqref{eq:red-folded} with primitive support and palindromic
coefficients. By hypothesis its data is of known type, so by
\autoref{lem:unfold}(iv) the data of \eqref{eq:familya} is of known type, and
\autoref{lem:red-members} finishes the proof.
\end{proof}

The remaining sections therefore work with a primitive palindromic support
$S$, coefficients $c_j=c_{k-j}>0$ on $S$, and either value of $\theta$, and
show that global periodicity forces known type.
